\documentclass[11pt,openany]{book}

\usepackage[T1]{fontenc}
\usepackage[utf8]{inputenc}
\usepackage[a4paper,margin=28mm]{geometry}
\usepackage{setspace}
\usepackage{amsmath,amssymb}
\usepackage{array}
\usepackage{longtable}
\usepackage[hidelinks]{hyperref}

\onehalfspacing
\newcolumntype{L}[1]{>{\raggedright\arraybackslash}p{#1}}

\newcommand{\Mten}{M_{10}}
\newcommand{\Yfour}{Y_4}
\newcommand{\Xsix}{X_6}
\newcommand{\Pcoh}{P_{\mathrm{coh}}}
\newcommand{\Qinc}{Q_{\mathrm{inc}}}
\newcommand{\Hq}{\mathcal H_{\mathrm{q79}}}
\newcommand{\WW}{Q_{\mathrm{WW}}}
\newcommand{\Atrace}{\mathcal A_0}
\newcommand{\Lshared}{L_{\mathrm{shared}}}
\newcommand{\Soff}{\mathcal S_{\mathrm{off}}}
\newcommand{\status}[1]{\textsc{#1}}

\hypersetup{
  pdftitle={The Book on Modal Triplet Theory: A Typed Interpretive Overview},
  pdfauthor={Peter Nero}
}

\title{\Huge The Book on Modal Triplet Theory\\[2mm]
\Large A Typed Interpretive Overview of the Current Program\\[4mm]
\large What is proved, reconstructed, calibrated, and still open}
\author{Peter Nero}
\date{July 2026\\Version 10}

\begin{document}
\frontmatter
\maketitle

\chapter*{Abstract and Current Status}
\addcontentsline{toc}{chapter}{Abstract and Current Status}

Modal Triplet Theory (MTT) is a proposed architecture for spectral reduction,
coherent-sector stability, and comparison among physical encodings.  Its
rigorous core is a dimension-neutral Hilbert-bundle framework with compatible
vertical operators, a joint spectral projector, separately stated fixed-point
and equilibrium gates, controlled complementary-mode estimates, and typed
descent and recovery maps.  A canonical physical realization uses a supplied
ten-dimensional bundle $\Mten\to\Yfour$ with compact six-dimensional fiber
$\Xsix$ and a Lorentzian four-dimensional base.

The current geometry program combines a local rank-three comparison field
$\WW\in\Gamma(\operatorname{Hom}(TP,TI))$ with a selected q79 degree-three
spectral-cover carrier.  Both exhibit a $1+2+3$ rank profile, but their global
bundle-and-connection intertwiner remains to be proved.  The q79 branch has an
exact trace split and a local binary-dihedral spin lift; strict global Spin
closure is still an obstruction calculation.

Computationally, MTT closes embedded renormalized-Standard-Model equivalence at
the declared one-shared-physical-primitive/profile standard.  This includes the
finite matrix architecture, charged-Yukawa profile, CKM, electroweak,
threshold, finite-Dirac, and multi-loop precision packets.  It is not strict
zero-knob selection of the observed universe.  Standard BRST/Faddeev--Popov
quantization and measured renormalized profile data are admitted at this
closure tier.

This book is interpretive, but every interpretation is labeled.  It does not
claim that MTT has derived the Born rule, Einstein gravity, a UV-complete
quantum gravity, the origin of time, three families, cosmology, or all of
string theory from projection alone.

\chapter*{Revision Note for This Edition}
\addcontentsline{toc}{chapter}{Revision Note for This Edition}

\begin{description}
\item[Supersedes.] \emph{The Book on Modal Triplet Theory}, version~9
(January 2026).

\item[Reason.] Version~9 treated triplet minimality, ten-dimensional geometry,
particle and force assignments, basin/Born probabilities, Gaussian quantum
gravity, framework containment, and cosmological scenarios as more established
than the underlying papers allowed.  It also used the retired few-TeV execution
chain, conflated the shared phase circle with time-like bookkeeping, and
predated the selected q79, finite-algebra, HYM, SM-profile, and foundational-
geometry successor theorems.

\item[Resolution.] Version~10 is rebuilt around the corrected Foundation,
Fixed Points I--VI, Projection--Admissibility, Signature, Scale, Relationship
Atlas, Theta/Execution, q79, proto-spinor, closure-strain, action, and current
SM ledgers.  It uses explicit claim tiers throughout, corrects the
$3+3$ versus $3\times3$ distinction, replaces literal circle--lens--nil
topology by typed carrier roles, and reports numerical closure at its declared
profile standard.

\item[Retained result.] The central interpretive idea survives: one can study
stable effective descriptions through shared spectral, geometric, and
fixed-point structures, and several nontrivial MTT realizations now have exact
or reproducible certificates.

\item[Remaining boundary.] The leading unresolved promotions are the strict
global q79 Spin certificate, the world-in-world/strain-to-q79 same-source
intertwiner, a selected normalized action source, strict no-knob and unique-
branch selection, and independent MTT derivations of quantum probability,
gauge quantization, gravity, and cosmology.
\end{description}

\noindent\textbf{Major withdrawals from version 9.}
The following are not carried forward as results: universal proof that three
filters or three families are inevitable; identification of compact phase with
physical time; literal $S^1\times\mathrm{Lens}\times\mathrm{Nil}$
compactification; Born probabilities from basin volume alone; instantaneous
bilocal microcausality; all-loop finite and unitary quantum gravity from a
Gaussian filter; a derived Big-Bang replacement; and the old $4.2$--$5$ TeV
cross-sector numerical chain.

\tableofcontents

\mainmatter

\part{How to Read the Current MTT Program}

\chapter{Ambition with a Type System}

MTT remains ambitious.  It asks whether stable effective physics can be
organized by one mathematical language of vertical spectra, coherent
projection, admissible continuation, and controlled reduction.  The ambition
does not authorize every attractive interpretation.  The first discipline of
the current program is therefore a type system for claims.

\section{Seven claim types}

\begin{description}
\item[Assumption.] Data inserted into a realization: a manifold, metric,
operator, representation, measure, action, or boundary condition.
\item[Conditional theorem.] A proved consequence of explicitly listed
hypotheses on a stated domain.
\item[Characterization.] A necessary form for any object satisfying the
premises, without proving existence or selection.
\item[Reconstruction.] Recovery of a known framework after target-compatible
data are supplied.
\item[Controlled reduction.] A map from one description to another with a
domain and error estimate.
\item[Calibration or profile equivalence.] Agreement using measured or
target-related data, including branch or scheme selection.
\item[Held-out prediction.] A value not used in construction, normalization,
scale choice, branch choice, or parameter fitting.
\end{description}

Interpretive correspondence is useful but is not an eighth theorem type.  It
is a proposed reading of already typed mathematics.

\section{The current one-paragraph verdict}

The internal spectral/control spine is rigorous under its hypotheses.  The
selected q79 and finite-algebra branches contain substantial exact structure.
The Standard Model is reconstructed at a declared profile-equivalence
standard with one shared physical primitive.  MTT has not yet shown that this
branch and these profile values are uniquely selected from a parameter-free
microscopic source.  It is therefore more than a metaphor and less than a
completed fundamental theory.

\section{Why this narrower statement is stronger}

A broad narrative can always absorb another analogy.  A typed program can
fail.  If an operator does not commute, a projector cannot be multiplied as
claimed.  If a relator closes with central sign $-1$, a strict Spin lift fails.
If normalized overlap rows do not match, the same-source bridge fails.  If a
registered held-out observable misses its uncertainty band, the numerical
branch fails.  The revised program gains credibility by creating such points
of contact.

\chapter{What Is Assumed}

No universal physical conclusion follows from ``coherence'' alone.  A concrete
MTT model declares at least:

\begin{enumerate}
\item the base, internal bundle, regularity class, and operator domains;
\item a strongly commuting family of vertical operators, or one total
internal operator, and its selected spectral cluster;
\item a well-posed stabilization flow and an invariant domain;
\item the relation, if any, between stabilization, physical time, and RG scale;
\item a physical action or evolution law, state space, observables, and
probability rule when physical claims are made; and
\item numerical provenance: sourced inputs, fitted inputs, branch selection,
uncertainty, and held-out outputs.
\end{enumerate}

These declarations are not blemishes.  Every physical theory has input
structure.  The real question is which inputs can later be related or selected
by a common theorem.

\part{The Rigorous Internal Spine}

\chapter{Hilbert Bundles, Vertical Operators, and Coherence}

\section{Dimension-neutral architecture}

Let $Y$ be a smooth base and $\mathcal E\to Y$ a real or complex Hilbert
bundle.  On a declared Sobolev space
\[
 \mathcal H=H^s(Y;\mathcal E),
\]
MTT uses three compatible vertical structures
\[
 (\mathcal E_i,A_i,P_i),\qquad i=1,2,3.
\]
The $A_i$ are nonnegative self-adjoint operators with specified domains.  If
they strongly commute, their spectral projectors may be multiplied:
\[
 \Pcoh=P_1P_2P_3,
 \qquad \Qinc=I-\Pcoh.
\]
Otherwise one must define a single total internal operator and take its
spectral projector.  Commutation is a hypothesis or theorem, never a gift of
notation.

The triplet is not automatically three coordinate manifolds.  Its entries may
be operators, bundle summands, representation lanes, or filtrations on one
internal carrier.

\section{Three parameters that must not be merged}

The corrected Foundation distinguishes:
\[
 \tau=\text{stabilization parameter},\qquad
 t=\text{physical time},\qquad
 \mu=\text{renormalization scale}.
\]
A heat or relaxation semigroup $R_\tau$ is not a Lorentzian propagator
$U(t_2,t_1)$, and neither is an RG trajectory.  A realization may relate them,
but it must provide the map, dimensions, domain, and error.

\section{Complementary-mode control}

On the noncoherent sector one seeks a semigroup estimate of the form
\[
 \|e^{\tau L_Q}\|\le M e^{-\gamma\tau},\qquad \gamma>0,
\]
or a suitable dissipativity bound.  This form remains meaningful for
nonnormal generators.  Eigenvalue signs alone do not control transient growth.
An internal gap controls internal complementary modes; it is not automatically
a four-dimensional momentum cutoff.

\chapter{Fixed Points Are Not All the Same}

Let $R_\tau$ be a well-posed stabilization map and define a projected step
\[
 T_\tau=\Pcoh R_\tau.
\]
A point satisfying $T_\tau\Psi_\ast=\Psi_\ast$ is a projected time-step fixed
point.  It is not automatically a steady solution of the upper flow.  A strict
Lyapunov identity or another stationarity argument is required to promote it
to an equilibrium.

Banach's theorem supplies existence and uniqueness only when $T_\tau$ is a
strict contraction on a nonempty complete invariant domain.  Schauder and
Darbo provide different existence routes and do not imply uniqueness.  This
is why the old slogan ``one inequality selects one universe'' has been
retired.  The current theory uses a ledger of independent gates.

\section{The six Fixed-Point papers}

\begin{longtable}{L{0.12\textwidth}L{0.73\textwidth}}
FP I & Geometry, Sobolev spaces, spectral projectors, semigroup well-posedness,
conditional existence and uniqueness routes, and Galerkin approximation.\\
FP II & Application to a supplied $4+6$ control realization, with a shared
circle counted as bundle data and equilibrium promotion kept separate.\\
FP III & Deterministic disturbance amplitudes, stochastic noise powers,
noncoherent damping floors, exact OU statements, and conditional homogenization.\\
FP IV & Curved spectral-cluster persistence, $Q\mathcal RP$ leakage,
intrinsic Karcher centroids, first-order modulation, and work-limited exits.\\
FP V & Frozen linear covariance, canonical-correlation estimates, declared
admissibility margins, and Gaussian exit bounds.\\
FP VI & A typed synthesis separating inherited theorem, conditional physical
completion, and interpretation.\\
\end{longtable}

These results constitute the strongest general mathematical achievement of
MTT.  They do not select a particle spectrum or fundamental action by
themselves.

\chapter{Projection, Descent, and Recovery}

Suppose $P_0:X\to Y_0$ is an initial reduction, $P_t:X\to Y_t$ a final
reduction, and $\Phi_t:X\to X$ an upper evolution.  Four questions must be
separated:

\begin{enumerate}
\item Does a right section choose one compatible upper representative?
\item Does a left decoder recover the actual upper state?
\item Does an autonomous reduced map $F_t:Y_0\to Y_t$ exist?
\item Is that reduced map injective, or do effective labels merge?
\end{enumerate}

Noninjectivity of $P_0$ does not forbid a right section.  Autonomous descent
exists exactly when $P_t\Phi_t$ is constant on the equivalence classes of
$P_0$.  Probability requires a measure and disintegration; entropy requires a
state or measure plus a dynamical or coarse-graining law.  Neither appears
from noninjectivity alone.

\part{Geometry: Local Carrier and Global Candidate}

\chapter{The Canonical \texorpdfstring{$4+6$}{4+6} Realization}

The canonical physical specialization is
\[
 \pi:\Mten\longrightarrow\Yfour,
 \qquad \dim\Yfour=4,
 \qquad \dim X_x=6.
\]
The physical completion supplies a globally hyperbolic Lorentzian base and a
compact Riemannian internal fiber.  This is a selected realization, not a
dimension theorem of the abstract Hilbert-bundle framework.

If the internal space factorizes, the dimensions of its factors must add to
six.  A shared phase circle is represented by a principal $U(1)$ bundle or
Hermitian line bundle.  It is not appended as a seventh coordinate.

\section{Signature comes from a principal symbol}

A Gram tensor
\[
 G_{\mu\nu}=\langle D_\mu\Psi,D_\nu\Psi\rangle
\]
is positive semidefinite and cannot be a Lorentzian metric.  Causal signature
is read from the principal symbol of a selected hyperbolic physical equation,
schematically
\[
 \sigma_{\mathrm{pr}}(x,\xi)
 =g^{\mu\nu}(x)\xi_\mu\xi_\nu I+\text{gauge blocks}.
\]
The current signature theorem establishes compatibility and stability of one-
time hyperbolicity.  It does not derive three spatial dimensions.

\chapter{World-in-World Comparison Geometry}

The lawful version of the world-in-world idea is a field between two oriented
rank-three bundles:
\[
 \WW\in\Gamma\!\left(\operatorname{Hom}(TP,TI)\right).
\]
In local frames it is a $3\times3$ matrix with nine components.  This is not
manifold-dimension multiplication.  A rank-three bundle over a three-
dimensional base has six-dimensional total space.

At an invertible background, polar decomposition separates orientation and
strain:
\[
 \operatorname{Mat}(3,\mathbb R)
 =\mathfrak{so}(3)\oplus\operatorname{Sym}(3,\mathbb R),
 \qquad 9=3+6.
\]
Once an orthonormal flag is selected, the strain sector splits as
\[
 \operatorname{Sym}(3,\mathbb R)
 =\mathbb RI_3\oplus\mathcal D_0\oplus\Soff,
 \qquad 6=1+2+3.
\]
The scalar, traceless-diagonal, and symmetric off-diagonal projectors are
orthogonal for the Frobenius metric.  The $1+2+3$ split depends on the selected
flag; the fully $SO(3)$-invariant split is $1+5$.

Therefore
\[
 1+3\times3=(1+3)+(1+2+3)=4+6=10
\]
is an exact component identity after an ordering scalar is supplied.  It is
not a proof that these components are the tangent directions of $\Mten$.

\section{Iwasawa and the nil distinction}

For $GL^+(3,\mathbb R)$, Iwasawa decomposition has dimensions
\[
 \dim SO(3)=3,\qquad \dim A=3=1+2,\qquad \dim N=3.
\]
The upper-unitriangular $N$ is the three-dimensional Heisenberg nil group.
The symmetric off-diagonal strain space is also three-dimensional, but is not
itself that Lie algebra.  They are related by a background-dependent local
coordinate map.  This distinction prevents dimension matching from becoming
an unproved dynamical identity.

\chapter{The Selected q79 Carrier}

For the selected degree-three spectral cover $\pi:C\to B$, let
\[
 \mathcal A=\pi_*\mathcal O_C,
 \qquad \Atrace=\ker\operatorname{Tr}.
\]
The normalized trace projector gives
\[
 \mathcal A=\mathcal O\oplus\Atrace,
 \qquad \operatorname{rank}(\mathcal O,\Atrace,\mathcal A)=(1,2,3).
\]
The common carrier is
\[
 \Hq=\Lshared\otimes
 (\mathcal O\oplus\Atrace\oplus\mathcal A),
 \qquad \operatorname{rank}\Hq=6.
\]
This is a direct trace/trace-zero/full filtration.  A connected transitive
cover does not provide a global ordering of three individual sheets.

The q79 arithmetic branch includes the exact selected residue construction
leading to $q=79$ modulo the declared order-$448$ quotient.  Its value is a
branch theorem with stated assumptions, not a replacement for the geometry
selecting every physical observable.

\section{Local spin lift and global obstruction}

The signed sheet representation
\[
 \rho_+(\sigma)=\operatorname{sgn}(\sigma)P_\sigma
 \in SO(3)
\]
has inverse image $\operatorname{Dic}_3$ of order twelve in
$\operatorname{Spin}(3)$.  Exact quaternion lifts satisfy the local braid and
central-sign relations.  Strict global Spin closure still requires every
branch-complement relator to lift with the correct central sign, equivalently
the relevant obstruction class must vanish.  A shared-circle SpinC
cancellation is a different conclusion.

\section{Which global topology is current?}

The q79 Fu--Yau branch is the strongest selected physical compactification
candidate in the current corpus.  The space
$L(3,1)\times\mathrm{Nil}_3$ remains useful as an auxiliary or effective model.
They are not the same manifold; even their first Betti numbers differ.  The
literal product $S^1\times\mathrm{Lens}\times\mathrm{Nil}$ and literal
manifold nesting are retired as proof sources.

The shared circle is common $U(1)$ phase or holonomy data and is counted once.
Physical time is noncompact order in the Lorentzian completion.  A map
$\mathbb R\to S^1$ can synchronize a phase after its frequency is supplied,
but compact phase and time are not identical objects.

\chapter{The Missing Geometry Theorem}

The local real strain carrier and the complex q79 carrier share the rank
profile $1+2+3$.  Equality of the integer profile does not yet make them
objects of the same type.  One type-correct candidate is the complex-bundle
map
\[
 \mathfrak I_{\mathrm{WW}\to q79}:
 (\mathbb RI_3\oplus\mathcal D_0\oplus\Soff)
 \otimes_{\mathbb R}\mathbb C
 \longrightarrow
 \Lshared\otimes_{\mathbb C}
 (\mathcal O\oplus\Atrace\oplus\mathcal A).
\]
Equivalently, one could first select and prove a rank-six real form of the q79
carrier.  Either route is additional source data, not a consequence of the
rank count.  The resulting map must be a global bundle isomorphism and must
preserve:

\begin{enumerate}
\item transition functions and the selected flag;
\item the strain and HYM inner products;
\item covariant derivatives;
\item the Hessian and retarded operators; and
\item the normalized overlap kernels used by the numerical packets.
\end{enumerate}

This one theorem would connect the most attractive local interpretation to the
most developed global and numerical branch.  Repeating the rank count does not
advance it.

\part{From Carrier to Physical Models}

\chapter{Proto-Spinors Without Automatic Particles}

If oriented rank-three frame data must retain the nontrivial class of
$\pi_1(SO(3))\cong\mathbb Z_2$, an ordinary $SO(3)$ vector description loses
that loop memory.  The minimal connected cover retaining it is
\[
 \operatorname{Spin}(3)\cong SU(2).
\]
This is a conditional double-cover theorem.  It does not prove that nature
must begin with rank three, and it does not yet create a Lorentzian fermion.

A physical spacetime spinor requires a Lorentzian spin or SpinC structure,
tetrad, compatible connection, hyperbolic Dirac operator, statistics, state
space, and observables.  Weyl, Dirac, Majorana, and twistor descriptions are
conditional chart languages once their chirality, pairing, reality, and
conformal assumptions are supplied.

Circle, lens, and nil now name carrier roles:

\begin{description}
\item[Circle.] Shared phase or holonomy.
\item[Lens.] Finite, projective, or signed-sheet transport.
\item[Nil.] Triangular transport, anchoring, or termination in a selected
operator filtration.
\end{description}

This taxonomy need not be a literal global product topology.

\chapter{Closure Strain, Scalars, and Mass}

The exact local achievement is the strain normal form.  A positive closure
Hessian proves local restoring cost on its declared slice.  It does not prove
that one direction is the Higgs or that an eigenvalue is a physical mass.

A pole mass comes from a canonically normalized quadratic action.  If
\[
 S_4^{(2)}=\frac12\int\!\sqrt{|g_4|}\,
 \left(Z_{AB}\partial_\mu\varphi^A\partial^\mu\varphi^B
 -M^2_{AB}\varphi^A\varphi^B\right),
\]
then the tree-level squared masses are eigenvalues of
$Z^{-1/2}M^2Z^{-1/2}$.  A closure Hessian becomes $M^2$ only through a
normalized source map.

The selected finite-algebra computation begins with three rank-four real
scalar-doublet modules.  A q79/proto-spinor alignment projector selects one
rank-four module and removes eight unwanted real scalar directions at the
declared profile tier.  This is the relevant one-Higgs selection result.  It
does not follow from Hessian positivity alone.

\chapter{A Regime-Local Action, Not a Derived Final Action}

The corrected action paper supplies a ten-dimensional EFT ansatz on a declared
domain.  A representative two-derivative form contains
\[
 \begin{aligned}
 \int_{\Mten}\sqrt{|g_{10}|}\,\bigg[
 &\frac{R_{10}}{2\kappa_{10}^2}
 -\frac14 k_{ab}F^a_{MN}F^{bMN}
 -\frac12G_{AB}D_M\phi^A D^M\phi^B-V(\phi)\\
 &{}+i\bar\Psi\Gamma^M D_M\Psi
 -\bar\Psi\mathcal M(\phi)\Psi
 \bigg]
 \end{aligned}
\]
plus flux, dilaton, gauge-fixing, ghost, boundary, and higher-order terms.

The metric and Einstein--Hilbert term are inputs.  Curvature is defined by a
connection; $\nabla S\ne0$ alone does not imply curvature or Frobenius failure.
Discrete internal spectra require compact resolvent or an equivalent spectral
compactness theorem.  Consistent four-dimensional truncation requires the
discarded equations to vanish on the retained ansatz, or an explicit
Schur--Feshbach error bound.

\part{The Standard Model Result at Its Actual Tier}

\chapter{Finite Algebra and Matrix Architecture}

The selected finite branch uses
\[
 \mathcal A_F=\mathbb C\oplus\mathbb H\oplus M_3(\mathbb C)
\]
on an explicit three-family particle--antiparticle carrier.  In the declared
counting, the current finite construction is a 96-dimensional complex finite
Hilbert space with a $96\times96$ profile Dirac operator satisfying
self-adjointness, grading, KO-dimension-six reality, order-zero, and structural
order-one checks.

The $27\times27$ qutrit--Weyl/minimal matrix is a finite response architecture,
not a standard universal size demanded by quantum mechanics.  Its dimensions
come from the selected tensor and sector organization.  Exact algebraic and
profile checks establish that the construction can carry the required gauge,
family, flavor, phase, Higgs, and threshold data at the adopted standard.

The anomaly analysis selects the usual anomaly-free hypercharge direction on
the chosen spectrum and preserves the familiar quotient
\[
 (SU(3)\times SU(2)\times U(1)_Y)/\mathbb Z_6.
\]
This is a theorem about the selected finite representation.  The local
$1+2+3$ strain dimensions alone do not derive these representations.

\chapter{What ``True SM Closure'' Means Here}

The current non-looping ledger reports 12/12 obligations closed for
\emph{embedded renormalized-Standard-Model equivalence at the
one-shared-physical-primitive/profile standard}.  The locked baseline includes:

\begin{itemize}
\item the $27\times27$ finite matrix architecture;
\item charged Yukawa magnitudes at SM-parity/profile tier;
\item selected CKM weight rows and their uncertainty certificate;
\item electroweak and direct Higgs/threshold rows;
\item the internal threshold ledger;
\item the finite Dirac operator at profile tier;
\item selected multi-loop mass-scheme and threshold transport;
\item a positive-definite precision covariance object; and
\item a five-arrow renormalized-SM observable functor.
\end{itemize}

This is a serious reconstruction result.  It shows that the selected MTT
branch can encode and reproduce the renormalized SM data consistently through
one finite operator architecture and common transport policy.

It is not a strict no-knob theorem.  Measured renormalized profile parameters
are admitted downstream, and one shared physical primitive fixes the declared
normalization.  Standard SM BRST/Faddeev--Popov quantization is imported rather
than derived from MTT.

\section{Parameters and knobs}

Three standards must remain separate:

\begin{description}
\item[Profile parity.] Can the architecture represent the accepted SM data
and reproduce them under the frozen conventions?
\item[One-shared-primitive equivalence.] Can the sectors be tied consistently
with one common physical normalization while measured profile values are
admitted?
\item[Strict no-knob selection.] Does one microscopic MTT source select the
branch, dimensionful normalization, and all values without using those data?
\end{description}

The first two are closed at the stated tier.  The third is an upgrade target.
Calling all three ``closure'' without the qualifier caused much of the earlier
confusion.

\section{Neutral sector and strong CP}

The neutral sector has a selected profile realization with declared holonomy
and absolute-scale inputs.  Strict selection of the phase, scale, and
Dirac/Majorana ontology remains open.  The strong-CP program has advanced
substantially through the q79 heterotic axion and anomaly analysis, but the
strict fully selected numerical nonperturbative source is not yet closed.

\chapter{Theta Closure and the End of the Few-TeV Chain}

The corrected Theta program transports measured source data with a common
multi-loop scheme to $Q=M_t$.  In its stated normalization the gauge-profile
targets are
\[
 \frac{I_2}{I_1}=0.5110273\pm0.0001231,
 \qquad
 \frac{I_3}{I_1}=0.158335\pm0.001098.
\]
They are calibrated profile coordinates, not first-principles predictions.
The overlaps are weighted gauge-kinetic coefficients; unit normalization of
representatives does not force all of them to one.

An auxiliary round-$S^2$ and Heisenberg-nil ansatz can be retargeted to these
ratios and passes its declared dimensionless spectral-floor test.  The twistor
route reproduces the same quadratic functionals under shared conventions and
is therefore a conditional cross-check, not independent normalization.

The old $4.2$--$5$ TeV crossing is withdrawn.  It is not a coherence scale,
internal gap, quantum-gravity scale, or prediction source.  The auxiliary
volume coefficient $20.07064R_1^3$ belongs only to its stated
$S^1\times S^2\times\mathrm{Nil}_3$ ansatz and is not the q79 Fu--Yau volume.

The weak-angle relation derived from the same gauge ratio is an exact
round-trip identity, not a held-out prediction.  A genuine test requires an
MTT source for the ratio that does not use electroweak gauge data.

\part{Quantum, Field-Theoretic, and Gravitational Interpretation}

\chapter{Quantum Mechanics: What Is and Is Not Reconstructed}

Projection, coherent basins, and stable sectors can represent alternatives,
records, and reduced states.  This gives useful language for decoherence and
measurement models.  It does not produce quantum probability by itself.

A full quantum reconstruction must supply:

\begin{enumerate}
\item a complex Hilbert space and observables;
\item unitary dynamics or a specified open-system law;
\item canonical commutation relations where relevant;
\item measurement instruments and update rules; and
\item the Born functional or a theorem equating basin weights with trace
probabilities for all relevant states and measurements.
\end{enumerate}

Basin volume depends on a measure.  Counting basins is not probability.
Noninjective projection is not uncertainty.  A deterministic branch narrative
must also confront Bell's theorem: MTT cannot be advertised as a Bell-local
hidden-variable completion.  A consistent global account would need to state
its nonfactorization, contextuality, boundary-value, retrocausal, or other
mechanism explicitly while preserving no-signaling.

The phrase ``no fundamental dice'' is therefore a philosophical preference in
this book, not a theorem of current MTT.

\chapter{QFT and the Renormalized SM Functor}

The strongest QFT result is a reconstruction on the selected branch.  The
finite representation and action data feed standard renormalized-SM
quantization and multi-loop transport, producing a typed observable functor.
This is enough for profile equivalence.  It does not derive the path integral,
BRST complex, or continuum interacting QFT from projection.

Fiberwise coherent compression can preserve locality of a supplied upper local
net when the coherent sector is invariant and the reduction is genuinely
fiberwise.  This is a meaningful AQFT-style descent theorem.  A spatially
nonlocal instantaneous kernel is not microcausal merely because it leaves the
principal symbol unchanged.  A local mediator is the consistent route to
causal overlap dynamics.

\chapter{Worldsheets, Strings, and Fu--Yau Geometry}

The proto-spinor/worldsheet paper now defines a $C^2$ bridge map between
declared configuration spaces.  If its derivative $L$ intertwines the two
Hessians,
\[
 L^*H_{\mathrm{ws}}L=H_{\mathrm{ps}},
\]
then the quadratic diagnostics agree up to a cubic remainder near alignment.
This is a local conditional bridge.

It does not prove that two-dimensionality is inevitable or that a complete
string theory follows.  Worldsheet Weyl invariance, beta functions, central
charge, ghosts, modular invariance, spin-structure sums, anomalies, tadpoles,
unitarity, and target-space matching remain independent gates.

Fu--Yau/Hull--Strominger mathematics supplies a legitimate class of heterotic
flux backgrounds and is the right established setting for the selected q79
branch.  MTT currently reconstructs and selects data inside this setting; it
does not contain all of string or M-theory as a proved theorem.

\chapter{General Relativity and Gravity}

The canonical completion supplies a Lorentzian base and a gravitational action
or field equation.  Coherent reduction may then produce a four-dimensional
effective Einstein sector if discarded equations, moduli, warping, gauge
vectors, and error terms are controlled.

The internal exact-branch transverse-traceless support packet is a real
structural result.  Physical Newton/Planck normalization and the complete
stress-energy response remain separate.  Gravity is therefore a conditional
physical completion of the MTT spine, not something projection has already
derived.

The image of gravity as ``elastic response'' can remain an interpretation of
the selected action and strain fields.  It is not a substitute for deriving
the Einstein equations and their observables.

\chapter{Quantum Gravity: Present Boundary}

Filtered propagators and proper-time factors can soften selected Euclidean or
TT integrals.  This does not prove all-loop finiteness, unitarity, constructive
existence, asymptotic completeness, BRST/BV consistency, causal support, or a
physical graviton sector.

The old book's ``ontic Gaussian damping'' and completed UV-finite quantum-
gravity language are withdrawn.  Current QG papers are conditional models and
diagnostic bridges.  They may inspire a completion, but they are not evidence
that MTT already exceeds established quantum-gravity programs.

\chapter{Time, Entropy, and Cosmology}

Physical time is the noncompact parameter of a selected Lorentzian evolution.
Stabilization order $\tau$, RG order, chart refinement, and compact circle
phase do not become time merely by being ordered.

An arrow of time requires an asymmetric law, state, or boundary condition.
A Lyapunov function can orient stabilization on its domain; entropy requires a
measure or state and a coarse-graining or thermodynamic theorem.  These may be
related in a concrete model, but are not currently identical.

The Great Coherence, pre-Big-Bang continuation, inflation as relaxation, dark
matter as modal solitons, baryogenesis from geometry, and cosmic renewal are
interpretive scenarios.  None is a current MTT cosmological prediction.  A
physical cosmology needs a selected action, solution, initial-state measure,
perturbation theory, transport to observables, and comparison with data.

\part{Relationships to Established Frameworks}

\chapter{A Typed Atlas, Not a Universal Superset}

\begin{longtable}{L{0.19\textwidth}L{0.25\textwidth}L{0.42\textwidth}}
Framework & Current relation & Missing promotion\\\hline
\endfirsthead
Framework & Current relation & Missing promotion\\\hline
\endhead
General Relativity & Conditional physical completion and reduction target &
Selected action, consistent truncation, stress response, and observables.\\
Quantum Mechanics & Interpretive reconstruction bridge & Born functional,
measurement instruments, and full equivalence.\\
QFT/AQFT & Conditional reconstruction with locality descent & Continuum
interacting theory, states, renormalization, scattering, and anomalies.\\
Standard Model & Embedded renormalized-SM equivalence at profile standard &
Strict no-knob source and unique observed-branch selection.\\
Effective Field Theory & Schur--Feshbach and mode-reduction template &
Nonlinear error, symmetry, power counting, matching, and RG control.\\
Kaluza--Klein theory & Conditional internal-mode reconstruction & Selected
compactification, warping, mixing, and stabilized moduli.\\
Noncommutative geometry & Selected finite-algebra realization & Full spectral-
triple axioms and unique geometric source where not already certified.\\
String/Calabi--Yau & Encoding and reconstruction program & Complete worldsheet
consistency, dualities, branes/fluxes, and low-energy matching.\\
Hull--Strominger and Fu--Yau & Conditional solution-slice reconstruction & Unique
selected background and same-source operator/value emission.\\
Quantum gravity & Open conditional program & Unitarity, causality, continuum
existence, BRST/BV, and physical observables.\\
LQG, causal sets, asymptotic safety & Interpretive or conditional bridges &
Explicit controlled maps and bidirectional observable equivalence.\\
\end{longtable}

Multiple MTT realizations reproducing the same target data create
underdetermination.  Agreement proves realizability, not inevitability.  A
canonical source theorem must select among inequivalent realizations.

\part{What Has Been Achieved and What Comes Next}

\chapter{The Genuine Achievements}

\section{Methodological}

MTT now has a disciplined separation of assumptions, conditional theorems,
reconstructions, profiles, and predictions.  This sounds modest, but it
prevents dozens of attractive category errors and makes the program auditable.

\section{Functional analytic}

The corrected Foundation and Fixed Points series provide a reusable framework
for vertical spectral reduction, nonnormal-safe damping, multiple existence
and uniqueness routes, curved-projector persistence, leakage, covariance, and
admissibility exits.

\section{Geometric}

The local $3\times3$ comparison carrier and its exact $3+6$ and $1+2+3$
decompositions are verified.  The q79 degree-three cover supplies a separate
global $1+2+3$ trace carrier, exact local signed-sheet spin data, and a
developed Fu--Yau/HYM branch.  The missing bridge is now finite and explicit.

\section{Computational}

The selected finite-algebra and $27\times27$ architecture pass extensive exact
and numerical checks.  Embedded renormalized-SM equivalence closes at the
declared one-shared-primitive/profile standard, including multi-loop precision
transport and covariance.  This is unusual in breadth for a speculative
framework, even though it is not yet no-knob prediction.

\section{Epistemic}

Several former claims were disproved or retired: the right-inverse argument,
Lorentzian Gram metric, extra-circle dimension count, few-TeV chain, independent
weak-angle test, literal Lens--Nil/q79 identity, and automatic Born/QG claims.
Finding and correcting these failures is part of the achievement.

\chapter{The Decisive Frontier}

The next work should follow this order.

\section{Global q79 spin decision}

Lift a presentation of the q79 branch complement and compute the central sign
of every relator.  This decides strict Spin closure.  If the obstruction does
not vanish, state the SpinC alternative without relabeling it as Spin.

\section{World-in-world to q79 intertwiner}

Construct $\mathfrak I_{\mathrm{WW}\to q79}$ with transition, metric,
connection, Hessian, retarded, and overlap compatibility.  This is the most
important foundational bridge for the numerical program.

\section{One selected action source}

Use one q79/Fu--Yau action, gauge fixing, regulator, scale, and scheme to emit
normalized zero modes and all required finite response rows.  Verify the
discarded equations or a reduction error bound.

\section{Strict value and branch selection}

Freeze the source before looking at the held-out data.  Count continuous and
discrete inputs honestly.  Then test future or unused common-scheme observables
with uncertainty and a preregistered acceptance rule.

\section{Probability and quantization}

Construct the measure/state and prove, or falsify, a basin--trace equality.
Derive the BRST and path-integral structure from MTT only if one common source map
actually produces it.

\section{Gravity and cosmology}

After the action is selected, derive the Lorentzian equations, constraints,
stress response, Planck normalization, cosmological backgrounds, perturbations,
and observables.  Interpretive cosmology comes last, not first.

\chapter{Falsifiability Without Promotional Predictions}

MTT can fail before new experimental data arrive:

\begin{itemize}
\item the q79 relator signs may obstruct the desired strict Spin lift;
\item no global connection-preserving $1+2+3$ intertwiner may exist;
\item the selected action may fail consistent truncation or anomaly control;
\item normalized source rows may miss the locked profile;
\item parameter counting may collapse back to the SM input count;
\item the basin--trace or locality bridge may be impossible; or
\item independent held-out observables may disagree with the frozen branch.
\end{itemize}

A genuine empirical prediction must list every datum used in geometry,
normalization, scheme choice, branch choice, and parameter fitting.  Agreement
with a value used anywhere in that chain is a profile check or postdiction.

\chapter{An Interpretive Picture That Survives}

The surviving picture is quieter than version~9 and more interesting.
Physical descriptions may be stable projections of richer geometric and
operator data.  Different theories may be charts or reconstructions on
overlapping domains without one being literally contained in another.
Triplet structure can be a selected rank filtration rather than a mystical
universal number.  The shared circle can coordinate phase without becoming
time.  Particles can be studied as stable excitations without pretending that
stability alone determines their quantum numbers.  Gravity can be compared to
strain response after an action is supplied.  Measurement can be modeled as a
physical interaction without announcing a derivation of probability.

This interpretation is useful because each sentence points to a theorem that
could be built.  It asks for maps instead of resemblance, sources instead of
fitted rows, and observables instead of slogans.

\chapter{Afterword}

MTT is no longer best described as ``one field, three filters, one inequality,
one inevitable world.''  The current program is richer:

\begin{quote}
One typed spectral architecture; several exact internal modules; one selected
q79/Fu--Yau and finite-algebra branch; a closed renormalized-SM profile
equivalence; and a short list of decisive bridges still to prove.
\end{quote}

That statement gives up inevitability but gains contact with mathematics.  It
also preserves the original impulse of the book: physics should become more
coherent as its hidden assumptions are exposed.  The next success will not be
another sweeping interpretation.  It will be one of the missing maps actually
constructed and independently checked.

\appendix

\chapter{Current Claim Ledger}

\begin{longtable}{L{0.41\textwidth}L{0.20\textwidth}L{0.27\textwidth}}
Claim & Current status & Controlling boundary\\\hline
\endfirsthead
Claim & Current status & Controlling boundary\\\hline
\endhead
Hilbert-bundle spectral/control spine & Conditional theorem & Operator domains,
commutation, gaps, invariant domains.\\
Projected fixed points & Conditional theorem & Existence route and fixed-point
type must be stated.\\
Equilibrium promotion & Conditional theorem & Strict Lyapunov or equivalent
stationarity gate.\\
$4+6$ physical realization & Canonical input & Not derived from abstract MTT.\\
$3\times3=3+6$ local split & Proved & Nonsingular comparison background.\\
$1+2+3$ strain split & Proved with flag & Flag is selected input/source.\\
q79 $1+2+3$ trace carrier & Selected theorem & Degree-three cover.\\
Local $\operatorname{Dic}_3$ lift & Proved & Strict global relators still open.\\
Fu--Yau/HYM branch & Selected realization & Rank-two literal continuum witness;
stronger transfer/uniqueness guarded.\\
Finite matrix and Dirac operator & Closed at profile tier & Selected finite
representation and profile rows.\\
Embedded renormalized SM & Closed at one-shared-primitive/profile standard &
Measured profile and imported SM quantization admitted.\\
Strict no-knob SM & Open upgrade & Unique source and branch selection.\\
Born rule & Open & Measure/state and basin--trace theorem.\\
Einstein gravity & Conditional completion & Selected Lorentzian action and
controlled reduction.\\
UV-complete quantum gravity & Open & Full constructive, causal, gauge, and
unitarity obligations.\\
Cosmological origin scenario & Interpretation & Selected solution and
observable pipeline absent.\\
\end{longtable}

\chapter{Notation and Guardrails}

\begin{description}
\item[$\Mten\to\Yfour$.] Canonical supplied physical bundle with compact
six-dimensional fiber $\Xsix$.
\item[$\WW$.] Rank-three world-in-world comparison field; never the
noncoherent projector.
\item[$\Pcoh$.] Joint coherent spectral projector.
\item[$\Qinc$.] Complementary noncoherent projector $I-\Pcoh$.
\item[$\tau,t,\mu$.] Stabilization, physical time, and RG scale; distinct by
default.
\item[$\Lshared$.] Common phase/holonomy line; not an extra dimension or time.
\item[$\Soff$.] Real symmetric off-diagonal strain subspace; distinct from
the structure sheaf $\mathcal O$ used in the q79 carrier.
\item[$\mathcal A,\Atrace$.] Rank-three q79 cover algebra and its rank-two
trace-zero summand.
\item[$\Hq$.] Shared q79 rank-six $1+2+3$ carrier.
\item[Profile standard.] A reconstruction tier admitting measured
renormalized coordinates and one shared physical primitive.
\item[Strict standard.] A source-selected, no-target-input theorem with a
frozen branch and held-out validation.
\end{description}

\backmatter

\chapter{References}

\begin{thebibliography}{99}

\bibitem{Foundation2026}
P.~Nero, \emph{Modal Triplet Theory: Foundations}, version~7, 2026.

\bibitem{FixedPoints2026}
P.~Nero, \emph{Fixed Points I--VI}, corrected editions, 2026.

\bibitem{Projection2026}
P.~Nero, \emph{The Projection--Admissibility Principle: Descent, Recovery,
and Structural Constraints}, version~2, 2026.

\bibitem{Signature2026}
P.~Nero, \emph{Lorentzian Base Compatibility and Signature Stability in the
MTT Fixed-Point Realization}, version~2, 2026.

\bibitem{Atlas2026}
P.~Nero, \emph{Modal Triplet Theory: A Typed Relationship Atlas}, version~3,
2026.

\bibitem{GeometryReconciliation2026}
P.~Nero, \emph{MTT Foundational Geometry Reconciliation}, theorem and
verification packet, 2026.

\bibitem{ProtoSpinor2026}
P.~Nero, \emph{The Proto-Spinor: Conditional Spinorial Closure and the q79
Interface}, version~5, 2026.

\bibitem{WorldInWorld2026}
P.~Nero, \emph{World-in-World Genesis: Local Comparison Geometry and a
Globalization Program}, version~5, 2026.

\bibitem{ClosureStrain2026}
P.~Nero, \emph{Closure-Strain Geometry: Local Normal Forms and Conditional
Matter Encodings}, version~6, 2026.

\bibitem{Action2026}
P.~Nero, \emph{Closure Geometry and a Regime-Local Ten-Dimensional Action
Ansatz}, version~4, 2026.

\bibitem{SMClosure2026}
P.~Nero, \emph{MTT Current True SM Closure Consolidated Ledger}, theorem and
verification packet, 2026.

\bibitem{Theta2026}
P.~Nero, \emph{Theta Closure and Execution I--V}, corrected editions, 2026.

\bibitem{FuYau}
J.-X.~Fu and S.-T.~Yau, ``The theory of superstring with flux on non-Kahler
manifolds and the complex Monge--Ampere equation,''
\emph{Journal of Differential Geometry} \textbf{78} (2008), 369--428.

\bibitem{LawsonMichelsohn}
H.~B.~Lawson and M.-L.~Michelsohn, \emph{Spin Geometry}, Princeton University
Press, 1989.

\bibitem{Connes}
A.~Connes, \emph{Noncommutative Geometry}, Academic Press, 1994.

\bibitem{Polchinski}
J.~Polchinski, \emph{String Theory}, Volumes 1--2, Cambridge University Press,
1998.

\bibitem{WeinbergQFT}
S.~Weinberg, \emph{The Quantum Theory of Fields}, Cambridge University Press,
1995--2000.

\end{thebibliography}

\end{document}
